\documentclass[10pt,a4paper]{amsart}
\usepackage[margin=19mm]{geometry}
\usepackage{amsmath,amssymb,mathtools}
\usepackage[scr=rsfs,cal=euler]{mathalfa}
\usepackage{xfrac}
\usepackage[unicode,hidelinks,bookmarks=false]{hyperref}
\newcommand{\ie}{i.e.\ }
\newcommand{\cf}{cf.\ }
\newcommand{\resp}{respectively\ }
\newcommand{\Subsection}{\subsection}
\providecommand{\nobreakdash}{\nobreak\hbox{-}}
\setlength{\emergencystretch}{2em}
\multlinegap0pt
\usepackage{mathtools}
\usepackage{amssymb}
\usepackage{tikz}
\newtheorem{theorem}{Theorem}
\newcommand{\N}{\mathbb{N}}
\newcommand{\R}{\mathbb{R}}
\title{\bfseries Henstock--Kurzweil Gauge Integral in the Non-Gaussian Regime:\\ A Machine-Verified Construction}
\author{Yuri N. Berdinsky}
\date{Source: arXiv:2609.10793v1 (CC BY 4.0)}
\begin{document}
\maketitle

\noindent\emph{Source and licence.} \bfseries Henstock--Kurzweil Gauge Integral in the Non-Gaussian Regime:\\ A Machine-Verified Construction. Version arXiv:2609.10793v1 (CC BY 4.0), distributed by arXiv under the Creative Commons Attribution 4.0 International licence, http://creativecommons.org/licenses/by/4.0/, as recorded in the arXiv licence metadata for this exact version and shown on the abstract page https://arxiv.org/abs/2609.10793v1. Full licence notice: \texttt{licenses/arxiv-2609.10793-CC-BY-4.0.txt}.\par\smallskip

\noindent\textit{Editorial scope.} Counted unit: the article's theorem thm:diff (source0.tex lines 527--547). The article's complete original proof of it is delivered with it (source0.tex lines 549--581, one \texttt{proof} environment of 33 lines). No mathematics is written, repaired or re-derived: the statement and the proof are the article's own lines, in the article's own notation, and no retained line is substituted or re-flowed.  Retained uncounted as the source-local context the proof needs: lines 470--481.  Nothing was omitted from any retained span, so the package carries no omission record.  No retained line contains a citation, so the wrapper carries no bibliography and no bibliography entry.  No retained line contains a figure environment, an image-inclusion command or drawing code, so no image is delivered and the package declares no asset dependency.  Editorial formatting only, in addition: an AMS \texttt{amsart} preamble that restates the article's own declarations and macros used by the retained lines, namely the environments \texttt{theorem} on separate counters, together with the standard packages those lines need.  The retained environments print in this document as Theorem 1 in order of appearance, whereas the article prints the same items with the numbering of its own full document.  The complete native source of the article as distributed in the arXiv e-print for this exact version is bundled unchanged as \texttt{sources/209950/source0.tex}.
\begin{theorem}[Existence and finiteness]
\label{thm:finite}
Let $\omega>0$, $j\in\R$, $\lambda\ge0$ and $k\in\N$. Then the function
$\varphi\mapsto\varphi^{k}e^{-\frac{\omega}{2}\varphi^{2}+j\varphi-\lambda\varphi^{4}}$ is
absolutely integrable. In particular $I(\omega,j,\lambda)$ is finite, and so are all the
moments
\begin{equation}
\label{eq:moments}
M_m(\omega,j,\lambda)=\int_\R \varphi^{m}\,
e^{-\frac{\omega}{2}\varphi^{2}+j\varphi-\lambda\varphi^{4}}\,d\varphi .
\end{equation}
\end{theorem}

\begin{theorem}[Differentiation under the integral sign]
\label{thm:diff}
Let $\omega>0$, $j\in\R$.
\begin{enumerate}
\item For every $\lambda>0$ and every $m\in\N$,
\[
\frac{d}{d\lambda}\,M_m(\omega,j,\lambda)=-\,M_{m+4}(\omega,j,\lambda).
\]
\item Consequently $\lambda\mapsto M_m(\omega,j,\lambda)$ is of class $C^{\infty}$ on
$(0,\infty)$ and
\[
\frac{d^{k}}{d\lambda^{k}}\,I(\omega,j,\lambda)=(-1)^{k}M_{4k}(\omega,j,\lambda).
\]
\item At the boundary point $\lambda=0$ the one-sided derivative exists and
\begin{equation}
\label{eq:derzero}
\partial_\lambda I(\omega,j,0)
=-\int_\R \varphi^{4}\,e^{-\frac{\omega}{2}\varphi^{2}+j\varphi}\,d\varphi .
\end{equation}
\end{enumerate}
\end{theorem}

\begin{proof}
(1) Fix $\lambda>0$ and work on the ball $|\lambda'-\lambda|<\lambda/2$, on which
$\lambda'>\lambda/2$. The $\lambda'$-derivative of the integrand is
$-\varphi^{m+4}e^{-\frac{\omega}{2}\varphi^{2}+j\varphi-\lambda'\varphi^{4}}$, and its
absolute value is dominated, uniformly in $\lambda'$ in that ball, by
$(1+\varphi^{2m})\varphi^{4}e^{-\frac{\omega}{2}\varphi^{2}+j\varphi-\frac{\lambda}{2}\varphi^{4}}$,
which is integrable by Theorem~\ref{thm:finite}. The dominated-convergence theorem for
parametric integrals applies and gives the formula.

(2) Induction on the order: the derivative of $M_m$ is $-M_{m+4}$, which is itself
differentiable by (1); an induction on $n$ shows $M_m\in C^{n}((0,\infty))$ for all $n$.

(3) The interior argument fails at $\lambda=0$ because a two-sided neighbourhood of $0$
contains negative couplings, for which the integral diverges. Instead we estimate the
difference quotient directly. Write $w_\lambda(\varphi)$ for the integrand, so that
$w_\lambda=e^{-\lambda\varphi^{4}}w_0$. For $u\ge0$ one has the elementary inequalities
\[
0\;\le\;e^{-u}-1+u\;\le\;u^{2},
\]
the first from $e^{-u}\ge 1-u$, the second from $e^{u}\ge 1+u$ and
$(1+u)^{-1}\le 1-u+u^{2}$. With $u=\lambda\varphi^{4}$ this gives, for $\lambda>0$,
\[
\Big|\frac{w_\lambda(\varphi)-w_0(\varphi)}{\lambda}+\varphi^{4}w_0(\varphi)\Big|
=\frac{e^{-u}-1+u}{\lambda}\,w_0(\varphi)
\le \lambda\,\varphi^{8}w_0(\varphi),
\]
whence
\[
\Big|\frac{I(\omega,j,\lambda)-I(\omega,j,0)}{\lambda}+M_4(\omega,j,0)\Big|
\le \lambda\,M_8(\omega,j,0)\xrightarrow[\lambda\downarrow0]{}0 .
\]
This is exactly \eqref{eq:derzero}.
\end{proof}

\end{document}
